\honest{Heat vs. Motion}{Eliezer Yudkowsky}{2008}

After yesterday's post about anger, here is a simpler example of reduction crossing an apparent difference in kind: heat and motion. ``Heat is motion'' seems obvious now, but the kinetic theory was once controversial, against the belief in a caloric fluid, and earlier still the theory of heat was phlogiston. \nb{Right in outline. John Herapath's and John James Waterston's kinetic theories of 1820 and 1843 both failed peer review.}

Sadi Carnot, working within the caloric theory, found the most efficient possible heat engine in 1824. Someone who knows heat and mechanics as separate subjects, as Carnot did, feels that heat and motion are independent properties of matter.

A physicist from the future says: ``Where there is heat, there is motion, and vice versa.'' You might read this as a bridging law, by which motion creates caloric and caloric pushes on things. \nb{The caloric theorists themselves held that caloric could be neither created nor destroyed, so they could not have read it this way.} Or you might suppose that heat and motion are the same thing. \nb{Francis Bacon separated these two readings in 1620: ``it must not be thought that heat generates motion, or motion heat \ldots\ but that the very essence of heat \ldots\ is motion and nothing else.''}

Thinker 1 says the two words have different meanings. Thinker 1 confuses the quotation with the referent: ``2 + 2'' and ``4'' are different strings with the same value, and the morning star and the evening star can be one planet, which you find out with a telescope, not by examining your beliefs. Thinker 2 says he can imagine a world where heat and motion come apart, so they are different properties. He mistakes a rule of inference for a law of physics, and then takes what is conceivable for what is logically possible. To Carnot a world where heat and motion come apart was conceivable; to Feynman it would make one's head explode. \nb{By 1832 Carnot himself had written that heat ``is a movement among the particles of bodies.'' Kripke had analysed this very example in 1972, explaining why heat seems as if it might not have been molecular motion. The post does not mention Kripke.}

In fairness to philosophers, some have said this. I quote Hilary Putnam: ``Conceivability is no proof of logical possibility.'' But ``water'' seems to me to be used in two senses in Putnam's two paragraphs, and Putnam ``seems to be saying'' that after the experiments, water is ``automatically redefined'' to mean H2O; one could hold instead that it means ``whatever is really in that bottle next to me.'' \nb{Putnam's paper says that ``water'' referred to H2O in 1750, before anyone knew it, and that the meaning did not change with the discovery. The bottle reading is close to Putnam's own account, in which water is whatever is the same liquid as ``this liquid'' in the glass.} ``I believe the above has already been said as well?'' \nb{It had. The two readings correspond to the distinction that two-dimensional semantics makes between considering a world as actual and as counterfactual.}

One thing can feel like two when you do not know the reduction. Could you prove heat is motion if the physicists took it back tomorrow? Saying ``Maybe heat and motion are the same thing!'' is easy; explaining how is hard. ``Reductionism is easy, reduction is hard.''
